\chapter{INTERFACES DO COMPILADOR (3 MODALIDADES)}

Conforme especificado no enunciado oficial do projeto MiniPar 2026.2, o sistema oferece três modalidades distintas de interface: Linha de Comando (CLI), Primeira Interface Web (Desenvolvedor) e Segunda Interface Web (Didática).

\section{Modalidade 1: Interface de Linha de Comando (CLI) e Menu Terminal}

A interface por linha de comando (\texttt{main.py} e \texttt{menu.py}) é orientada à automação em esteiras de integração contínua (CI) e à execução ágil de testes diretamente pelo console:
\begin{itemize}
    \item Modo Direto: permite compilar qualquer arquivo \texttt{.minipar} passando seu caminho como argumento:
    \begin{lstlisting}[language=bash]
python main.py tests/oficiais/teste07_quicksort_iterativo_tac.minipar
    \end{lstlisting}
    \item Menu Interativo de Testes: executado via \texttt{python menu.py}, exibe uma lista numerada dos 11 testes oficiais, permitindo executar todos os testes em lote ou disparar um teste individual com captura de entrada do usuário.
    \item Coordenação Multi-Terminal: o módulo \texttt{test\_orchestrator.py} coordena múltiplos terminais abertos para cenários distribuídos, monitorando o progresso de cada participante através de estado compartilhado em \texttt{output/.test\_session.json}.
\end{itemize}

\section{Modalidade 2: Primeira Interface Web (Desenvolvedor)}

A primeira interface Web é voltada ao perfil de desenvolvedor e pode ser acessada na rota:
\[ \texttt{http://127.0.0.1:5000/ide} \]

Essa interface apresenta:
\begin{itemize}
    \item Editor de código baseado no CodeMirror com realce de sintaxe;
    \item Painel lateral com lista de arquivos de testes;
    \item Abas superiores com visualização simultânea do código C++ gerado e da saída em Três Endereços (TAC);
    \item Terminal inferior integrado para exibição da saída padrão (\texttt{stdout}) e mensagens de erro do compilador.
\end{itemize}

A Figura \ref{fig:ide_dev} ilustra a interface de desenvolvedor, permitindo editar programas MiniPar, compilar com exibição lado a lado de C++ e TAC, e executar com saída no terminal integrado.

\begin{figure}[htbp]
    \centering
    \framebox[\textwidth][c]{%
        \parbox{0.9\textwidth}{%
            \centering\vspace{1.5cm}
            \textcolor{red}{<Inserir print atualizado da tela da Primeira Interface Web (IDE do Desenvolvedor) no navegador em http://127.0.0.1:5000/ide>}\\
            \vspace{0.3cm}
            \textit{\small Nota: Capturar a tela do navegador acessando a rota /ide com um programa carregado, exibindo as abas de C++ e TAC e o terminal.}
            \vspace{1.5cm}
        }
    }
    \caption{Primeira Interface Web (IDE do Desenvolvedor -- \texttt{/ide})}
    \label{fig:ide_dev}
\end{figure}

\section{Modalidade 3: Segunda Interface Web (Didática / Educacional)}

O compilador inclui a Segunda Interface Web Didática, concebida para atender ao aprendizado prático de compiladores. Ela está disponível na rota raiz e didática do servidor:
\[ \texttt{http://127.0.0.1:5000/}\quad\text{e}\quad\texttt{http://127.0.0.1:5000/didatico} \]

\subsection{Inspiração Visual e Filosofia de Design}
Inspirada no compilador educacional do Instituto Federal de Santa Catarina (IFSC), a interface apresenta um layout limpo, moderno, responsivo e modular, com painéis ajustáveis e controles organizados por fase.

\subsection{Controles Didáticos Passo a Passo (Botões de Fase)}
Diferentemente de compilações monolíticas em que todo o processo ocorre em caixa-preta, a interface didática decompõe o processo de compilação em botões dedicados que podem ser acionados passo a passo:

\begin{enumerate}
    \item Tokens -- Lexer (Botão): apresenta uma tabela interativa contendo todos os tokens identificados pelo analisador léxico, detalhando a categoria sintática (\texttt{TokenType}), o lexema original, a linha e a coluna exatas, além de ferramentas de busca e filtros por categoria.
    \item AST -- Árvore Sintática (Botão): em vez de apenas um dump textual ou ASCII, a interface renderiza uma árvore sintática interativa em SVG/D3.js. Os nós possuem cores temáticas por categoria (comandos, expressões, classes), são colapsáveis com suporte a expansão sob clique, navegação por arrasto (\textit{pan}) e ampliação dinâmica (\textit{zoom}). Árvores extensas iniciam recolhidas para manter a tela limpa, com botão \textit{"Expandir Tudo"} para inspeção completa.
    \item Análise Semântica (Botão): exibe o status da validação semântica, lista de diagnósticos com linha/coluna e a Tabela de Símbolos completa, discriminando escopos léxicos, nomes, tipos e assinaturas.
    \item TAC -- Código de 3 Endereços (Botão): exibe o código intermediário linear destrinchado instrução por instrução com destaque sintático para registradores temporários e rótulos de controle.
    \item Código C++ Gerado (Botão): exibe o código-fonte C++20 transpilado gerado em memória pelo componente \texttt{ICodegen}.
    \item Compilar e Baixar Executável (Botão): dispara a compilação nativa isolada no backend via \texttt{g++ -O3} e disponibiliza o download direto do binário executável (\texttt{.exe} no Windows ou executável ELF no Linux).
    \item Executar (Botão): executa o binário compilado em um contêiner de execução no servidor, com suporte à passagem de dados pré-digitados no painel de \textit{Terminal \& Execução} (uma resposta por linha, preservando espaços e vírgulas).
\end{enumerate}

\begin{figure}[htbp]
    \centering
    \framebox[\textwidth][c]{%
        \parbox{0.9\textwidth}{%
            \centering\vspace{2cm}
            \textcolor{red}{<Inserir print da tela da Segunda Interface Web Didática com a Árvore Sintática (AST) interativa expandida e botões de fase didáticos no navegador>}\\
            \vspace{0.3cm}
            \textit{\small Nota: Capturar a tela acessando http://127.0.0.1:5000/ com um dos exemplos carregados (ex: Teste 02 ou Teste 03) exibindo a árvore em D3.js.}
            \vspace{2cm}
        }
    }
    \caption{Segunda Interface Web Didática do MiniPar 2026.2 (Visualizador Gráfico Interativo de AST)}
    \label{fig:ide_didatica_ast}
\end{figure}
